\documentclass[10pt,a4paper]{article}

% --------------------------------------------------
% CONFIGURACIÓN GENERAL
% --------------------------------------------------
\usepackage[utf8]{inputenc}
\usepackage[T1]{fontenc}
\usepackage[spanish]{babel}
\usepackage{lmodern}
\usepackage{microtype}
\usepackage[
    left=1.40cm,
    right=1.40cm,
    top=1.05cm,
    bottom=1.05cm
]{geometry}

\usepackage{xcolor}
\usepackage{enumitem}
\usepackage{titlesec}
\usepackage[hidelinks]{hyperref}

\pagestyle{empty}
\setlength{\parindent}{0pt}
\setlength{\parskip}{0pt}
\renewcommand{\familydefault}{\sfdefault}

% Paleta de colores
\definecolor{principal}{HTML}{243E53}
\definecolor{texto}{HTML}{292D32}
\definecolor{secundario}{HTML}{626B73}
\definecolor{linea}{HTML}{CBD3DA}

% --------------------------------------------------
% TÍTULOS DE SECCIÓN
% --------------------------------------------------
\titleformat{\section}
    {\fontsize{10}{12}\selectfont\bfseries\color{principal}}
    {}
    {0pt}
    {}
    [\vspace{1pt}{\color{linea}\titlerule[0.45pt]}]

\titlespacing*{\section}{0pt}{8pt}{4pt}

% --------------------------------------------------
% LISTAS
% --------------------------------------------------
\setlist[itemize]{
    leftmargin=1.25em,
    label=\textcolor{principal}{\textbullet},
    labelsep=0.55em,
    itemsep=1pt,
    topsep=2.5pt,
    parsep=0pt,
    partopsep=0pt
}

% --------------------------------------------------
% ELEMENTOS REUTILIZABLES
% --------------------------------------------------

% Empresa o institución y ubicación
\newcommand{\entidad}[2]{%
    \noindent
    \textbf{#1}\hfill
    {\small\color{secundario}#2}\par
}

% Cargo y fechas
\newcommand{\cargo}[2]{%
    \noindent
    \textit{#1}\hfill
    {\small\color{secundario}#2}\par
}

% Programa académico y fechas
\newcommand{\programa}[2]{%
    \noindent
    #1\hfill
    {\small\color{secundario}#2}\par
}

% Competencias y herramientas
\newcommand{\detalle}[2]{%
    \noindent
    \textbf{#1} #2\par
    \vspace{2pt}
}

\begin{document}

\fontsize{9.2}{11}\selectfont
\color{texto}

\begin{center}
    {\fontsize{23}{27}\selectfont
     \bfseries\color{principal}Valeria Soto Díaz\par}

    \vspace{4pt}

    {\fontsize{9.5}{12}\selectfont
     \bfseries Control de Gestión \textbar{} Retail \textbar{} Forecast de margen \par}

    \vspace{5pt}

    {\fontsize{8.5}{10.5}\selectfont\color{secundario}
     Santiago, Chile \textbullet{} +56 9 0000 0101 \textbullet{}
     \href{mailto:valeria.soto@example.com}{valeria.soto@example.com}\par}
\end{center}

\vspace{-3pt}
{\color{principal}\hrule height 1pt}
\vspace{2pt}

\section{RESUMEN PROFESIONAL}

Profesional de Ingeniería Comercial con experiencia en jefatura de Control de Gestión
en retail. Lidera el cierre y el forecast de margen, supervisa un equipo de tres
personas y coordina con Comercial la conciliación de diferencias de margen.

\section{COMPETENCIAS TÉCNICAS}

\detalle{Control de gestión}{Cierre, forecast de margen y control presupuestario.}

\section{HERRAMIENTAS TÉCNICAS}

\noindent\textbf{Excel avanzado}\par\vspace{2pt}
\detalle{Power BI}{Panel de cierre.}

\section{EXPERIENCIA PROFESIONAL}

\entidad{Retail Ejemplo}{}
\cargo{Jefa de Control de Gestión}{Ene 2020 -- Presente}

\begin{itemize}
    \item Lidera el cierre y el forecast de margen y supervisa un equipo de tres personas.
    \item Controla un presupuesto de CLP 8.000 millones.
    \item Redujo el plazo de cierre de 8 a 5 días con un panel en Power BI.
    \item Concilió diferencias de margen con Comercial y acordó un catálogo común.
    \item Corrigió un forecast de margen mediante alertas y validación de supuestos con Comercial.
\end{itemize}

\vspace{5pt}

\cargo{Analista de Finanzas}{Jul 2016 -- Dic 2019}

\section{FORMACIÓN ACADÉMICA}

\entidad{U. Ejemplo}{}
\programa{Ingeniería Comercial}{2011 -- 2015}

\end{document}
